\chapter*{Abstract}
\addcontentsline{toc}{chapter}{Abstract}

\begin{otherlanguage}{english}
At the beginning of 2024, a new research line was established within the Atmospheric Physics Group to investigate the possible existence of microbial life in liquid veins within ice \cite{mader2006}, \cite{rohde2007}, \cite{dani2012}. During its initial stages, the research focused on studying the movement of microorganisms found in stagnant water samples collected in the city of Córdoba, within a temperature range of 0 °C to 15 °C. The aim was to gain a better understanding of, among other aspects, their potential survival in environments close to the melting point of ice \cite{pedernera2024}.

Particular attention was given to the types of trajectories described by these microorganisms, their mean velocities, and the dependence of these quantities on the temperature of the medium. Their movement was monitored through video recordings obtained using a camera coupled to an optical microscope employed to observe the samples. Four types of microorganisms were identified across all the samples analyzed.

Currently, movement analysis is performed entirely manually using the freely available Tracker software. This process requires a considerable amount of time and limits the progress of the research. In this context, the development of an automated detection and tracking system is essential to optimize the processing of experimental data and enable the extraction of additional quantitative metrics, including velocity distributions, individual trajectories, acceleration variations, and changes in movement patterns in response to external stimuli, such as variations in the temperature of the medium.
\end{otherlanguage}
